% Ręcznie tworzona lista symboli - do użycia w wydruku pracy i na Overleafie
\flushleft

\textbf{\large Wykaz skrótów i akronimów}
\vspace{0.8em}

\begin{tabular}{ll}
    \textbf{AC} & \emph{Alternating Current} -- prąd przemienny \\
    \textbf{DC} & \emph{Direct Current} -- prąd stały \\
    \textbf{DT} & \emph{Digital Twin} -- Bliźniak Cyfrowy \\
    \textbf{ESR} & \emph{Equivalent Series Resistance} -- zastępcza rezystancja szeregowa \\
    \textbf{FS-MPC} & \emph{Finite-State Model Predictive Control} -- sterowanie predykcyjne ze skończonym zbiorem stanów \\
    \textbf{IAE} & \emph{Integral of Absolute Error} -- całka modułu błędu \\
    \textbf{ISE} & \emph{Integral of Squared Error} -- całka kwadratu błędu \\
    \textbf{ITAE} & \emph{Integral of Time-weighted Absolute Error} -- całka iloczynu czasu i modułu błędu \\
    \textbf{LPF} & \emph{Low-Pass Filter} -- filtr dolnoprzepustowy \\
    \textbf{MF-BB} & \emph{Model-Free Bang-Bang} -- bezmodelowe sterowanie przekaźnikowe \\
    \textbf{MSE} & \emph{Mean Squared Error} -- błąd średniokwadratowy \\
    \textbf{NMP} & \emph{Non-Minimum Phase} -- układ z dynamiką fazy nieminimalnej \\
    \textbf{OCP} & \emph{Over-Current Protection} -- zabezpieczenie nadprądowe dławika \\
    \textbf{OOP} & \emph{Object-Oriented Programming} -- programowanie obiektowe \\
    \textbf{OZE} & Odnawialne Źródła Energii \\
    \textbf{PI} & Regulator proporcjonalno-całkujący (\emph{Proportional-Integral}) \\
    \textbf{PSO} & \emph{Particle Swarm Optimization} -- Optymalizacja Rojem Cząstek \\
    \textbf{PWM} & \emph{Pulse Width Modulation} -- modulacja szerokości impulsów \\
    \textbf{RHPZ} & \emph{Right Half-Plane Zero} -- zero w prawej półpłaszczyźnie zmiennej zespolonej $s$ \\
    \textbf{Trend MF-BB} & Zmodyfikowany regulator MF-BB z filtrowanym trendem napięcia \\
    \textbf{ZOH} & \emph{Zero-Order Hold} -- ekstrapolator zerowego rzędu (opóźnienie pomiarowe 1 próbki) \\
\end{tabular}

\vspace{1.5em}
\textbf{\large Wykaz ważniejszych symboli i oznaczeń}
\vspace{0.8em}

\begin{tabular}{lll}
    \textbf{Symbol} & \textbf{Jedn.} & \textbf{Opis} \\
    \midrule
    $V_{\mathrm{in}}$ & V & Napięcie zasilania przekształtnika (strona dolnego napięcia) \\
    $L$ & H & Indukcyjność dławika magazynującego energię ($0{,}75~\mathrm{mH}$) \\
    $R_L$ & $\Omega$ & Pasożytnicza rezystancja uzwojenia dławika ($0{,}05~\Omega$) \\
    $C$ & F & Pojemność kondensatora filtru wyjściowego ($200~\mu\mathrm{F}$ nominalnie) \\
    $R_{\mathrm{load}}$ & $\Omega$ & Zastępcza rezystancja obciążenia wyjściowego ($47{,}619~\Omega$ dla $1{,}2~\mathrm{kW}$) \\
    $u_{\mathrm{dc}}(t)$ & V & Chwilowe napięcie na kondensatorze wyjściowym \\
    $i_L(t)$ & A & Chwilowy prąd płynący przez dławik \\
    $u_{\mathrm{ref}}$ & V & Wartość zadana napięcia wyjściowego ($240~\mathrm{V}$ / $160~\mathrm{V}$) \\
    $i_{\mathrm{des}}$ & A & Pożądany prąd cewki generowany w algorytmie MF-BB \\
    $s(t)$ & --- & Dwustanowy sygnał sterujący kluczami półprzewodnikowymi, $s \in \{0, 1\}$ \\
    $D$ & --- & Współczynnik wypełnienia impulsów (ang. \emph{duty cycle}) \\
    $w_i$ & --- & Współczynnik wagowy uchybu prądu w funkcji przełączającej \\
    $w_u$ & --- & Współczynnik wagowy tłumienia trendu napięcia wyjściowego \\
    $\Delta u_f$ & V & Dyskretny przyrost przefiltrowanego napięcia wyjściowego w kroku regulatora \\
    $f_{c,i}$ & Hz & Częstotliwość odcięcia filtru dolnoprzepustowego prądu ($321~\mathrm{Hz}$) \\
    $f_{c,u}$ & Hz & Częstotliwość odcięcia filtru dolnoprzepustowego napięcia ($321~\mathrm{Hz}$) \\
    $T_s$ & s & Okres próbkowania i wykonywania algorytmu sterowania ($10~\mu\mathrm{s}$) \\
    $\Delta t$ & s & Krok całkowania numerycznego fizyki obwodu ($0{,}5~\mu\mathrm{s}$) \\
    $I_{\max}, I_{\min}$ & A & Sprzętowe progi ogranicznika prądowego dławika ($\pm 20~\mathrm{A}$) \\
    $w$ & --- & Waga bezwładności w algorytmie PSO (ang. \emph{inertia weight}) \\
    $c_1, c_2$ & --- & Współczynniki uczenia kognitywnego i społecznego w roju PSO \\
    $\lambda_i$ & --- & Współczynnik kary za błąd prądu w funkcji celu Kroku 1 \\
    $w_{\mathrm{cross}}$ & --- & Współczynnik kary za zapad napięcia w funkcji celu Kroku 2 \\
\end{tabular}
